\documentclass[12pt]{report}
\usepackage[margin=1in]{geometry}
\usepackage{hyperref}
\title{Universal Hybrid Network Intrusion Detection System for Live, Offline, and Multi-Format Threat Analysis}
\author{Shaik}
\date{March 11, 2026}
\begin{document}
\begin{titlepage}
\centering
{\LARGE Universal Hybrid Network Intrusion Detection System for Live, Offline, and Multi-Format Threat Analysis \par}
\vspace{1cm}
{\large Author: Shaik \par}
{\large Institution: Institution to be updated \par}
{\large Supervisor: Supervisor to be updated \par}
{\large Project Version: research-snapshot-2026.03.11 \par}
\end{titlepage}
\chapter*{Abstract}
Universal NIDS is a repository-grounded hybrid intrusion-detection platform for live traffic, offline replay, artifact triage, machine learning, fusion-based alerting, and research-quality evidence retention.
\chapter{Introduction}
The project addresses heterogeneous evidence, operational constraints, and the need for a continuously updated research record.
\chapter{Literature Review}
The thesis draws on protocol-aware IDS, operational guidance, dataset-evaluation studies, online anomaly detection, and concept-drift literature \cite{paxson1998bro,scarfone2007idps,dreger2006dynamic,tavallaee2009kdd,sharafaldin2018cicids,ring2019survey,mirsky2018kitsune,seth2024conceptdrift}.
\chapter{System and Evaluation Summary}
Current training accuracy is 0.99536 with weighted F1 0.99560. Current evaluation accuracy is 0.99784 with weighted F1 0.99792. The implementation uses scapy>=2.5.0, PyYAML>=6.0.1, rich>=13.7.1, fastapi>=0.115.0, uvicorn[standard]>=0.30.6, joblib>=1.4.2, numpy>=1.26.4, pytest>=8.3.2, plotly>=5.24.1, pandas>=2.2.2, kaleido>=0.2.1, pypdf>=5.1.0, python-docx>=1.1.2, openpyxl>=3.1.5, beautifulsoup4>=4.12.3, lxml>=5.3.0, pefile>=2024.8.26, scikit-learn>=1.5.2, python-magic>=0.4.27, xgboost>=3.2.0.
\chapter{Conclusion}
Universal NIDS provides a practical hybrid architecture and a research-grade evidence trail while openly retaining unresolved gaps for future work.
\bibliographystyle{IEEEtran}
\bibliography{references}
\end{document}
